% Chapter 4, Section 1 _Linear Algebra_ Jim Hefferon
%  http://joshua.smcvt.edu/linearalgebra
%  2001-Jun-12
\section{幂零性}
本章将证明，每个方阵都相似于一个能写成两类简单矩阵之和的矩阵。
前一小节关注了第一类简单矩阵，即对角矩阵。
现在我们来考虑另一类。









\subsectionoptional{自复合}\index{self composition!of maps}\index{map!self composition}\index{transformation!composed with itself}\index{composition!self}
% \noindent\textit{This subsection is optional, although it is necessary
% for later material in this section and in the next one.}

由于线性变换 $\map{t}{V}{V}$ 的定义域与陪域相同，我们可以将
$t$ 与自身复合：
\( t^2=\composed{t}{t} \)，
以及 \( t^3=\composed{t}{\composed{t}{t}} \)，
等等。\appendrefs{function iteration}\spacefactor=1000 %
\begin{center}
  \includegraphics{jc/mp/ch5.1}
\end{center}
对变换的迭代使用上标幂记号 $t^j$，
与我们对其方阵表示所用的记号是相容的，
因为若 $\rep{t}{B,B}=T$，则 \( \rep{t^j}{B,B}=T^j \)。


\begin{example} \label{ex:DerivIter}
对于求导映射 \( \map{d/dx}{\polyspace_3}{\polyspace_3} \)，它由下式给出
\begin{equation*}
  a+bx+cx^2+dx^3\xmapsunder{d/dx} b+2cx+3dx^2
\end{equation*}
其二次幂就是二阶导数，
\begin{equation*}
  a+bx+cx^2+dx^3\xmapsunder{d^2/dx^2} 2c+6dx
\end{equation*}
其三次幂就是三阶导数，
\begin{equation*}
  a+bx+cx^2+dx^3\xmapsunder{d^3/dx^3} 6d
\end{equation*}
而任意更高次幂都是零映射。
\end{example}

\begin{example}
空间 $\matspace_{\nbyn{2}}$（即 $\nbyn{2}$ 矩阵构成的空间）中的这个变换
\begin{equation*}
  \begin{mat}
    a  &b  \\
    c  &d
  \end{mat}
  \mapsunder{t}
  \begin{mat}
    b  &a  \\
    d  &0
  \end{mat}
\end{equation*}
有如下二次幂
\begin{equation*}
  \begin{mat}
    a  &b  \\
    c  &d
  \end{mat}
  \mapsunder{t^2}
  \begin{mat}
    a  &b  \\
    0  &0
  \end{mat}
\end{equation*}
以及如下三次幂。
\begin{equation*}
  \begin{mat}
    a  &b  \\
    c  &d
  \end{mat}
  \mapsunder{t^3}
  \begin{mat}
    b  &a  \\
    0  &0
  \end{mat}
\end{equation*}
此后，$t^4=t^2$，$t^5=t^3$，依此类推。
\end{example}

\begin{example}
考虑平移变换 $\map{t}{\C^3}{\C^3}$。
\begin{equation*}
  \colvec{x \\ y \\ z} \mapsunder{t} \colvec{0 \\ x \\ y}
\end{equation*}
我们有
\begin{equation*}
  \colvec{x \\ y \\ z} \mapsunder{t} \colvec{0 \\ x \\ y}
                       \mapsunder{t} \colvec{0 \\ 0 \\ x}
                       \mapsunder{t} \colvec{0 \\ 0 \\ 0}
\end{equation*}
所以值域空间下降到平凡子空间。
\begin{equation*}
  \rangespace{t}=\set{\colvec{0 \\ a \\ b}\suchthat a,b\in\C}
  \qquad
  \rangespace{t^2}=\set{\colvec{0 \\ 0 \\ c}\suchthat c\in\C}
  \qquad
  \rangespace{t^3}=\set{\colvec{0 \\ 0 \\ 0}}
\end{equation*}
\end{example}

这些例子表明，经过若干次迭代之后，映射就趋于稳定。

\begin{lemma}  \label{le:RangeAndNullChains}
对任意变换 \( \map{t}{V}{V} \)，各次幂的值域空间构成一个递降链
\begin{equation*}
  V\supseteq \rangespace{t}\supseteq\rangespace{t^2}\supseteq\cdots
\end{equation*}
而零空间构成一个递升链。
\begin{equation*}
  \set{\vec{0}}\subseteq\nullspace{t}\subseteq\nullspace{t^2}\subseteq\cdots
\end{equation*}
此外，存在 \( k>0 \)，使得当幂小于 $k$ 时各包含关系都是真包含：~若
$j<k$，则 $\rangespace{t^j}\supset\rangespace{t^{j+1}}$
且
$\nullspace{t^j}\subset\nullspace{t^{j+1}}$；
而当
$j\geq k$ 时，则 $\rangespace{t^j}=\rangespace{t^{j+1}}$
且
$\nullspace{t^j}=\nullspace{t^{j+1}}$）。
\end{lemma}

\noindent （$k=1$ 的情形可能发生，例如当 $t$ 是恒等映射时，
此时两条链中没有一个包含关系是真包含。）

\begin{proof}
首先回忆，对任意映射，其值域空间的维数加上其零空间的维数等于
其定义域的维数。
于是，如果值域空间的维数收缩，那么零空间的维数必定上升。
这里我们只做值域空间这一半，其余部分留作
\nearbyexercise{exer:RangeAndNullChains}。

我们先证明
值域空间构成一条链。
若 $\vec{w}\in\rangespace{t^{j+1}}$，即有某个 $\vec{v}$ 使得
$\vec{w}=t^{j+1}(\vec{v})$，
则 $\vec{w}=t^{j}(\,t(\vec{v})\,)$。
于是 $\vec{w}\in\rangespace{t^{j}}$。

接下来我们验证“进一步”的性质：
链中的包含关系起初都是真包含，
而从某个幂 $k$ 开始值域空间彼此相等。
我们先证明：如果链中任意一对相邻的值域空间相等
\( \rangespace{t^{k}}=\rangespace{t^{k+1}} \)，
那么其后所有的值域空间也都相等：
\( \rangespace{t^{k+1}}=\rangespace{t^{k+2}} \)，等等。
这是因为
\( \map{t}{\rangespace{t^{k+1}}}{\rangespace{t^{k+2}}} \)
与
\( \map{t}{\rangespace{t^{k}}}{\rangespace{t^{k+1}}} \) 是定义域相同的同一个映射，
因此它有相同的值域
\( \rangespace{t^{k+1}}=\rangespace{t^{k+2}} \)
（由归纳法可知这对所有更高的幂都成立）。
因此，值域空间的链一旦停止严格递减，
从那一点起它就稳定了。

最后我们证明这条链必定最终停止递减。
每个值域空间都是前一个的子空间。
若要成为真子空间，它的维数就必须严格更低
（见 \nearbyexercise{exer:PropSubspStrictLowerDimen}）。
这些空间都是有限维的，所以这条链只能下降有限多步。
也就是说，幂 $k$ 至多是 $V$ 的维数。
\end{proof}

\begin{example}
$\polyspace_3$ 上的求导映射 $a+bx+cx^2+dx^3\xmapsunder{d/dx} b+2cx+3dx^2$
有如下值域空间链。
\begin{equation*}
  \rangespace{t^0}=\polyspace_3
    \:\supset\:\rangespace{t^1}=\polyspace_2
    \:\supset\:\rangespace{t^2}=\polyspace_1
    \:\supset\:\rangespace{t^3}=\polyspace_0
    \:\supset\:\rangespace{t^4}=\set{\zero}
    % =\rangespace{t^5}=\set{\zero}=
    % \,=\,\cdots
\end{equation*}
链中其后的各项都是平凡空间。
它有如下零空间链。
\begin{equation*}
  \nullspace{t^0}=\set{\zero}
  \:\subset\: \nullspace{t^1}=\polyspace_0
  \:\subset\:\nullspace{t^2}=\polyspace_1
  \:\subset\:\nullspace{t^3}=\polyspace_2
  \:\subset\:\nullspace{t^4}=\polyspace_3
\end{equation*}
其后的各项都是整个空间。
\end{example}

\begin{example} \label{exam:PolyRankFalls}
设 \( \map{t}{\polyspace_2}{\polyspace_2} \) 是映射
\( d_0+d_1x+d_2x^2 \mapsto 2d_0+d_2x. \)
如引理所描述的，在迭代下值域空间收缩
\begin{equation*}
  \rangespace{t^0}=\polyspace_2
    \quad
  \rangespace{t}=\set{a_0+a_1x\suchthat a_0,a_1\in\C}
    \quad
  \rangespace{t^2}=\set{a_0\suchthat a_0\in\C}
\end{equation*}
然后稳定下来，使得 $\rangespace{t^2}=\rangespace{t^3}=\cdots$。
零空间增长
\begin{equation*}
  \nullspace{t^0}=\set{0}
    \quad
  \nullspace{t}=\set{b_1x\suchthat b_1\in\C}
    \quad
  \nullspace{t^2}=\set{b_1x+b_2x^2\suchthat b_1,b_2\in\C}
\end{equation*}
然后稳定下来，$\nullspace{t^2}=\nullspace{t^3}=\cdots$。
\end{example}

\begin{example}
投影到前两个坐标的变换 \( \map{\pi}{\C^3}{\C^3} \)
\begin{equation*}
   \colvec{c_1 \\ c_2 \\ c_3}
     \mapsunder{\pi}
   \colvec{c_1 \\ c_2 \\ 0}
\end{equation*}
满足 \( \C^3\supset\rangespace{\pi}=\rangespace{\pi^2}=\cdots \)
且 \( \set{\zero}\subset\nullspace{\pi}=\nullspace{\pi^2}=\cdots\, \)，
其中的值域空间与零空间为
\begin{equation*}
  \rangespace{\pi}=\set{\colvec{a \\ b \\ 0}\suchthat a,b\in\C}
  \qquad
  \nullspace{\pi}=\set{\colvec{0 \\ 0 \\ c}\suchthat c\in\C}
\end{equation*}
\end{example}

\begin{definition}
设 \( t \) 是 \( n \) 维空间上的一个变换。
\definend{广义值域空间}\index{generalized range space}%
\index{range space!generalized}
（或 \definend{值域空间的闭包\/}\index{closure!of range space}%
\index{range space!closure of}）
是 $\genrangespace{t}=\rangespace{t^n}$。
\definend{广义零空间}\index{generalized null space}%
\index{null space!generalized}
（或 \definend{零空间的闭包\/}\index{closure!of null space}%
\index{null space!closure of}）
是 $\gennullspace{t}=\nullspace{t^n}$。
\end{definition}

这幅图给出了说明。
横轴给出变换的幂~$j$。
纵轴把
$t^j$ 的值域空间的维数表示为零点之上的距离，
因而也同时显示了零空间的维数，
因为两者之和等于定义域的维数 $n$。
\begin{center}
  % \includegraphics{jc/mp/ch5.2}
  \includegraphics{jc/mp/ch5.8}
\end{center}
在迭代下
秩下降而零度上升，
直到存在某个 $k$ 使得
映射达到稳定状态
$\rangespace{t^k}=\rangespace{t^{k+1}}=\genrangespace{t}$
且 $\nullspace{t^k}=\nullspace{t^{k+1}}=\gennullspace{t}$。
这最迟必定在第 $n$ 次迭代时发生。
% The steady state's distance above zero is the dimension of the
% generalized range space
% and its distance below $n$ is the dimension of
% the generalized null space.

\begin{exercises}
  \recommended \item
    给出零变换与恒等变换的值域空间链和零空间链。
    \begin{answer}
      对于零变换，无论空间是什么，值域空间链都是
      \( V\supset\set{\vec{0}}=\set{\vec{0}}=\cdots\, \)，
      而零空间链是 \( \set{\vec{0}}\subset V=V=\cdots\, \)。
      对于恒等变换，两条链分别是
      \( V=V=V=\cdots \) 与
      \( \set{\vec{0}}=\set{\vec{0}}=\cdots\, \)。
    \end{answer}
  \recommended\item
     对每个映射，
     给出值域空间链与零空间链，
     以及广义值域空间与
     广义零空间。
     \begin{exparts}
       \partsitem $\map{t_0}{\polyspace_2}{\polyspace_2}$,
         $a+bx+cx^2\mapsto b+cx^2$
       \partsitem $\map{t_1}{\Re^2}{\Re^2}$,
         \begin{equation*}
           \colvec{a \\ b}\mapsto\colvec{0 \\ a}
         \end{equation*}
       \partsitem $\map{t_2}{\polyspace_2}{\polyspace_2}$,
         $a+bx+cx^2\mapsto b+cx+ax^2$
       \partsitem $\map{t_3}{\Re^3}{\Re^3}$,
         \begin{equation*}
           \colvec{a \\ b \\ c}\mapsto\colvec{a \\ a \\ b}
         \end{equation*}
     \end{exparts}
     \begin{answer}
       \begin{exparts}
         \partsitem 将 $t_0$ 迭代两次
           $a+bx+cx^2\mapsto b+cx^2\mapsto cx^2$
           得到
           \begin{equation*}
             a+bx+cx^2\mapsunder{t_0^2}cx^2
           \end{equation*}
           而任何更高的迭代都是同一个映射。
           于是，虽然 $\rangespace{t_0}$ 是没有一次项的
           二次多项式构成的空间
           $\set{p+rx^2\suchthat p,r\in \C}$，
           且
           $\rangespace{t_0^2}$ 是纯二次多项式构成的空间
           $\set{rx^2\suchthat r\in \C}$，
           但链就在这里稳定下来：
           $\genrangespace{t_0}=\set{rx^2\suchthat r\in \C}$。

           至于零空间，
           $\nullspace{t_0}$ 是仅含一次项的
           多项式构成的空间 $\set{qx\suchthat q\in \C}$，
           而
           $\nullspace{t_0^2}$ 是一次多项式构成的空间
           $\set{p+qx\suchthat p,q\in \C}$。
           到此为止：$\gennullspace{t_0}=\nullspace{t_0^2}$。
         \partsitem 二次幂
           \begin{equation*}
             \colvec{a \\ b}
              \mapsunder{t_1}\colvec{0 \\ a}
              \mapsunder{t_1}\colvec{0 \\ 0}
           \end{equation*}
           是零映射。
           因此，值域空间链
           \begin{equation*}
             \Re^2
               \supset\set{\colvec{0 \\ p}\suchthat p\in\C}
               \supset\set{\zero}
               =\cdots
           \end{equation*}
           与零空间链
           \begin{equation*}
             \set{\zero}
               \subset\set{\colvec{q \\ 0}\suchthat q\in\C}
               \subset\Re^2
               =\cdots
           \end{equation*}
           的长度都是二。
           广义值域空间是平凡子空间，而
           广义零空间是整个空间。
         \partsitem 这个映射的迭代循环往复
           \begin{equation*}
             a+bx+cx^2
               \mapsunder{t_2} b+cx+ax^2
               \mapsunder{t_2} c+ax+bx^2
               \mapsunder{t_2} a+bx+cx^2
               \;\cdots
           \end{equation*}
           并且值域空间链与零空间链都是平凡的。
           \begin{equation*}
             \polyspace_2=\polyspace_2=\cdots
              \qquad
              \set{\zero}=\set{\zero}=\cdots
           \end{equation*}
           于是，显然，
           广义空间是 $\genrangespace{t_2}=\polyspace_2$
           与 $\gennullspace{t_2}=\set{\zero}$。
         \partsitem 我们有
           \begin{equation*}
             \colvec{a \\ b \\ c}
                \mapsto\colvec{a \\ a \\ b}
                \mapsto\colvec{a \\ a \\ a}
                \mapsto\colvec{a \\ a \\ a}
                \mapsto\cdots
           \end{equation*}
           于是值域空间链
           \begin{equation*}
             \Re^3
               \supset\set{\colvec{p \\ p \\ r}\suchthat p,r\in\C}
               \supset\set{\colvec{p \\ p \\ p}\suchthat p\in\C}
               =\cdots
           \end{equation*}
           以及零空间链
           \begin{equation*}
             \set{\zero}
                \subset\set{\colvec{0 \\ 0 \\ r}\suchthat r\in \C}
                \subset\set{\colvec{0 \\ q \\ r}\suchthat q,r\in \C}
                =\cdots
           \end{equation*}
           的长度都是二。
           广义空间就是上面每条链中最后所示的那个空间。
       \end{exparts}
      \end{answer}
  \item
    证明函数的复合满足结合律
    \( \composed{(\composed{t}{t})}{t}=\composed{t}{(\composed{t}{t})} \)，
    因而我们写 $t^3$ 时无须指明结合方式。
    \begin{answer}
      两者都把 \( x\mapsto t(t(t(x))) \)。
    \end{answer}
  \item \label{exer:PropSubspStrictLowerDimen}
    验证子空间的维数必定小于或等于
    其母空间的维数。
    验证若子空间是真子空间（子空间不等于
    母空间），则维数严格更小。
    \textit{（这用于
             \nearbylemma{le:RangeAndNullChains} 的证明中。）}
    \begin{answer}
      回忆一下，若 $W$ 是 $V$ 的子空间，则我们可以把
      $W$ 的任意一组基 $B_W$
      扩充为 $V$ 的一组基 $B_V$。
      由此第一句话立即可得。
      第二句话也不难：~$W$ 是 $B_W$ 的张成空间，
      若 $W$ 是真子空间，则 $V$ 不是 $B_W$ 的张成空间，
      所以 $B_V$ 至少比 $B_W$ 多一个向量。
    \end{answer}
  \recommended\item
    证明：若变换 $t$ 是非奇异的，则广义值域空间 $\genrangespace{t}$ 是
    整个空间，且广义零空间 $\gennullspace{t}$ 是平凡的。
    反过来“仅当”方向也成立吗？
    \begin{answer}
      “若”与“仅当”都成立。
      线性映射是非奇异的
      当且仅当它保持维数，也就是说，其值域的维数
      等于其定义域的维数。
      对于变换 $\map{t}{V}{V}$，这意味着
      该映射是非奇异的当且仅当它是满射：
      $\rangespace{t}=V$（从而 $\rangespace{t^2}=V$，等等）。
    \end{answer}
  \item \label{exer:RangeAndNullChains}
    验证 \nearbylemma{le:RangeAndNullChains} 中零空间的那一半。
    \begin{answer}
      零空间构成链，因为
      若 $\vec{v}\in\nullspace{t^j}$，则 $t^j(\vec{v})=\zero$，
      且 $t^{j+1}(\vec{v})=t(\,t^j(\vec{v})\,)=t(\zero)=\zero$，
      于是 $\vec{v}\in\nullspace{t^{j+1}}$。

      现在，
      零空间的“进一步”性质由它对值域空间成立这一事实
      结合前面的练习即可得出。
      由于 $\rangespace{t^j}$ 的维数加上
      $\nullspace{t^j}$ 的维数等于起始空间~$V$ 的维数~$n$，
      当值域空间的维数停止下降时，
      零空间的维数也不再变化。
      前面的练习表明，从这一点~$k$ 开始，
      链中的包含关系都不是真包含\Dash 零空间
      彼此相等。
    \end{answer}
  \recommended\item
    给出一个三维空间上的变换的例子，使其值域的维数为二。
    它的零空间是什么？
    对你的例子进行迭代，直到值域空间与零空间稳定。
    \begin{answer}
      （许多例子都正确，这里给出一个。）
      一个例子是实三元组上的平移算子
      \( (x,y,z)\mapsto (0,x,y) \)。
      零空间是所有以两个零开头的三元组。
      该映射在三次迭代后稳定。
     \end{answer}
  \item
      证明线性变换的值域空间与零空间未必不相交。
      它们有可能不相交吗？
      \begin{answer}
        求微分算子
        \( \map{d/dx}{\polyspace_1}{\polyspace_1} \) 的值域空间与零空间相同。
        要看一个二者不相交\Dash
        除零向量外\Dash 的例子，
        可考虑恒等映射，或任意非奇异映射。
      \end{answer}
\end{exercises}

















